% !PRINTED 234
\weberpage{234}
%<<CONT
\weq{10}{\begin{aligned}
\mu &\equiv \alpha = a\alpha' + c\beta'\\
\mu' &\equiv \beta = \partial\beta'
\end{aligned}\ (\mathrm{mod}\ n),}

worin $\beta'$ eine beliebige, zu $a$ teilerfremde Zahl bedeutet und
$\alpha'$ so bestimmt wird, daß $\alpha$, $\beta$ relativ prim werden, und
es ist auch leicht [nach §\,65, (3)] einzusehen, daß, so lange $a$,
$\partial$, $c$ dieselben bleiben, die nach (10) bestimmten Zahlenpaare
$\mu$, $\mu'$ derselben Reihe angehören.

Wir nennen die Zahlen $a$, $c$, $\partial$ (wie in §\,27) die
\gesperrt{Transformationszahlen} und $n$ den \gesperrt{Transformationsgrad}.

\gesperrt{Das Zahlensystem $a$, $c$, $\partial$ ist also vollständig charakteristisch
für eine Reihe}, und die Anzahl der Reihen ist gleich der Anzahl dieser
Zahlensysteme, woraus für die Zahl $\psi(n)$, (§\,63) folgt:

\weq{11}{\psi(n) = \Sigma \frac{a}{e}\varphi(e),}

wenn die Summe auf alle Divisoren $a$ von $n$ erstreckt wird.

Es ist von Interesse, diese Relation auch direkt zu beweisen, wobei die
Beschränkung auf ein ungerades $n$ wegfallen kann. Betrachten wir $\psi(n)$
jetzt als Zeichen für die Summe (11), so ergibt sich, wenn $n'$, $n''$
relativ prim sind, zunächst
\[
\psi(n')\psi(n'') = \psi(n'n''),
\]
und es bleibt also nur übrig, die Summe $\psi(n)$ für den Fall zu
bestimmen, daß $n=p^\pi$ eine Primzahlpotenz ist. In diesem Falle ist nun
$e$ gleich dem kleineren der beiden Divisoren $a$, $\partial$, und wenn
$a=\partial$ ist, $e=a$. Wenn wir also das erste und letzte Glied der Summe
$\psi(n)$ absondern, so erhalten wir
\[
\psi(p^\pi) = 1 + p^\pi + \sum_{1<a\leqq\sqrt{n}}\varphi(a) + \sum_{\sqrt{n}<a<n}\frac{a}{\partial}\varphi(\partial).
\]
Es ist aber
\[
\varphi(a) = a\frac{p-1}{p}, \qquad \varphi(\partial) = \partial\frac{p-1}{p},
\]
also
\[
\begin{aligned}
\psi(p^\pi) &= 1 + p^\pi + \frac{p-1}{p}\sum_{1<a<n} a\\
&= 1 + p^\pi + (p-1)(1+p+p^2+\cdots p^{\pi-2}) = p^\pi\Bigl(1+\frac{1}{p}\Bigr),
\end{aligned}
\]
%CONT>>
